%-----------------------------------------------------------------------------------------------------------------------------------------------%
%	The MIT License (MIT)
%
%	Copyright (c) 2021 Jitin Nair
%
%	Permission is hereby granted, free of charge, to any person obtaining a copy
%	of this software and associated documentation files (the "Software"), to deal
%	in the Software without restriction, including without limitation the rights
%	to use, copy, modify, merge, publish, distribute, sublicense, and/or sell
%	copies of the Software, and to permit persons to whom the Software is
%	furnished to do so, subject to the following conditions:
%	
%	THE SOFTWARE IS PROVIDED "AS IS", WITHOUT WARRANTY OF ANY KIND, EXPRESS OR
%	IMPLIED, INCLUDING BUT NOT LIMITED TO THE WARRANTIES OF MERCHANTABILITY,
%	FITNESS FOR A PARTICULAR PURPOSE AND NONINFRINGEMENT. IN NO EVENT SHALL THE
%	AUTHORS OR COPYRIGHT HOLDERS BE LIABLE FOR ANY CLAIM, DAMAGES OR OTHER
%	LIABILITY, WHETHER IN AN ACTION OF CONTRACT, TORT OR OTHERWISE, ARISING FROM,
%	OUT OF OR IN CONNECTION WITH THE SOFTWARE OR THE USE OR OTHER DEALINGS IN
%	THE SOFTWARE.
%	
%
%-----------------------------------------------------------------------------------------------------------------------------------------------%

%----------------------------------------------------------------------------------------
%	DOCUMENT DEFINITION
%----------------------------------------------------------------------------------------

% article class because we want to fully customize the page and not use a cv template
\documentclass[a4paper,12pt]{article}

%----------------------------------------------------------------------------------------
%	FONT
%----------------------------------------------------------------------------------------

% % fontspec allows you to use TTF/OTF fonts directly
% \usepackage{fontspec}
% \defaultfontfeatures{Ligatures=TeX}

% % modified for ShareLaTeX use
% \setmainfont[
% SmallCapsFont = Fontin-SmallCaps.otf,
% BoldFont = Fontin-Bold.otf,
% ItalicFont = Fontin-Italic.otf
% ]
% {Fontin.otf}

%----------------------------------------------------------------------------------------
%	PACKAGES
%----------------------------------------------------------------------------------------
\usepackage{url}
\usepackage{parskip} 	

%other packages for formatting
\RequirePackage{color}
\RequirePackage{graphicx}
\usepackage[usenames,dvipsnames]{xcolor}
\usepackage[scale=0.92,top=0.6in,bottom=0.6in]{geometry}

%tabularx environment
\usepackage{tabularx}

%for lists within experience section
\usepackage{enumitem}

% centered version of 'X' col. type
\newcolumntype{C}{>{\centering\arraybackslash}X} 

%to prevent spillover of tabular into next pages
\usepackage{supertabular}
\usepackage{tabularx}
\newlength{\fullcollw}
\setlength{\fullcollw}{0.47\textwidth}

%custom \section
\usepackage{titlesec}				
\usepackage{multicol}
\usepackage{multirow}

%CV Sections inspired by: 
%http://stefano.italians.nl/archives/26
\titleformat{\section}{\Large\scshape\raggedright}{}{0em}{}[\titlerule]
\titlespacing{\section}{0pt}{6pt}{5pt}

%for publications
\usepackage[style=authoryear,sorting=ynt, maxbibnames=2]{biblatex}

%Setup hyperref package, and colours for links
\usepackage[unicode, draft=false]{hyperref}
\definecolor{linkcolour}{rgb}{0,0.2,0.6}
\hypersetup{colorlinks,breaklinks,urlcolor=linkcolour,linkcolor=linkcolour}
\addbibresource{citations.bib}
\setlength\bibitemsep{1em}

%for social icons
\usepackage{fontawesome5}

%debug page outer frames
%\usepackage{showframe}


% job listing environments

\newenvironment{jobshort}[2]
    {
    \begin{tabularx}{\linewidth}{@{}l X r@{}}
    \textbf{#1} & \hfill &  #2 \\[3.75pt]
    \end{tabularx}
    }
    {
    }

% #1 = company, #2 = dates, #3 = job title (own bold-italic line, avoids wrap on long titles)
\newenvironment{joblong}[3]
    {
    \begin{tabularx}{\linewidth}{@{}l X r@{}}
    \textbf{#1} & \hfill &  #2 \\[1.5pt]
    \end{tabularx}
    \noindent\small\textbf{\textit{#3}}\par
    \vspace{1.5pt}
    \begin{minipage}[t]{\linewidth}
    \begin{itemize}[nosep,after=\strut, leftmargin=1em, itemsep=1.5pt,label=--]
    }
    {
    \end{itemize}
    \end{minipage}    
    }

%%% Работа с русским языком
\usepackage{cmap}					% поиск в PDF
\usepackage{mathtext} 				% русские буквы в формулах
\usepackage[T2A]{fontenc}			% кодировка
\usepackage[utf8]{inputenc}			% кодировка исходного текста
\usepackage[english,russian]{babel}	% локализация и переносы
\usepackage{indentfirst}
\frenchspacing


%----------------------------------------------------------------------------------------
%	BEGIN DOCUMENT
%----------------------------------------------------------------------------------------
\begin{document}

% non-numbered pages
\pagestyle{empty} 

%----------------------------------------------------------------------------------------
%	TITLE
%----------------------------------------------------------------------------------------

% \begin{tabularx}{\linewidth}{ @{}X X@{} }
% \huge{Your Name}\vspace{2pt} & \hfill \emoji{incoming-envelope} email@email.com \\
% \raisebox{-0.05\height}\faGithub\ username \ | \
% \raisebox{-0.00\height}\faLinkedin\ username \ | \ \raisebox{-0.05\height}\faGlobe \ mysite.com  & \hfill \emoji{calling} number
% \end{tabularx}

\begin{tabularx}{\linewidth}{@{} C @{}}
\Huge{Ногинов Иван Сергеевич} \\[7.5pt]
% \href{https://github.com/username}{\raisebox{-0.05\height}\faGithub\ username} \ $|$ \ 
% \href{https://linkedin.com/in/username}{\raisebox{-0.05\height}\faLinkedin\ username} \ $|$ \ 
% \href{https://mysite.com}{\raisebox{-0.05\height}\faGlobe \ mysite.com} \ $|$ \ 
\href{https://t.me/ExSchutZ}{\raisebox{-0.05\height}\faTelegram\ @ExSchutZ} \ $|$ \ 
\href{mailto:ex-schutz@yandex.ru}{\raisebox{-0.05\height}\faEnvelope \ ex-schutz@yandex.ru} \ $|$ \ 
\href{tel:+79211353765}{\raisebox{-0.05\height}\faMobile \ +7(921)135-37-65} \\
\end{tabularx}

%----------------------------------------------------------------------------------------
% EXPERIENCE SECTIONS
%----------------------------------------------------------------------------------------

%Interests/ Keywords/ Summary
\section{О себе}
Высококвалифицированный сетевой инженер с опытом обслуживания систем телемеханики и сетевой инфраструктуры в ПАО «Россети»: более 100 АРМ на базе ALT Linux в четырёх филиалах, внедрение ССПИ TOPAZ, проектирование SCADA-мониторинга СГП. Довожу решения от стадии проектирования до непосредственной реализации.

%----------------------------------------------------------------------------------------
%	SKILLS
%----------------------------------------------------------------------------------------
\section{Навыки}
\renewcommand{\arraystretch}{0.9}
\begin{tabularx}{\linewidth}{@{}>{\raggedright}p{3.4cm} X@{}}
Networking \& Protocols & \normalsize{Ethernet T568A/B, TCP/IP, VLAN/LAN (trunking, access ports), статическая маршрутизация, ACL, DNS, DHCP, NAT, SNMP, SSH, SSL/TLS, firewall, Wi-Fi, ModBUS/RTU, RS232/485, SNTP}\\[2pt]
Field \& Telemetry Equipment & \normalsize{Cisco (маршрутизаторы и коммутаторы), MOXA, ССПИ TOPAZ, УТМ «Телеканал-М2», УЦКС-1/2/1M, ИБП EATON APS3-321/APS6-321 (СГП)}\\[2pt]
Scripting & \normalsize{Python, Bash, CMD, PowerShell, SQL}\\[2pt]
OS \& Platforms & \normalsize{Linux, ALT Linux, Windows, NTFS, FAT32, ExFAT}\\[2pt]
Monitoring & \normalsize{MasterSCADA 3.9, OPC, Zabbix}\\[2pt]
% Databases & \normalsize{PostgreSQL 10+}\\[2pt]
Documentation \& Standards & \normalsize{Visio, Draw.io, формуляры приёмо-передачи данных, регламенты и инструкции по эксплуатации СГП и телемеханики}\\[2pt]
Languages & \normalsize{English B1}\\
\end{tabularx}
\renewcommand{\arraystretch}{1}

%----------------------------------------------------------------------------------------
%	EDUCATION
%----------------------------------------------------------------------------------------
\section{Образование \hfill \textnormal{\normalsize\href{https://disk.360.yandex.ru/client/aa/d_wbGuC3z0tBQeJg/}{Сертификаты. Курсы повышения квалификации.}}}
\begin{tabularx}{\linewidth}{@{}l l X@{}}	
2019 - 2023 & \textbf{НИУ МИЭТ} &  Кафедра ТКС, Инфокоммуникационные технологии и системы связи \hfill \normalsize \\ 
\end{tabularx}

%Experience
\section{Опыт работы}

\begin{joblong}{ПАО «Федеральная сетевая компания --- Россети»}{Ноябрь 2023 - ...}{Инженер по обслуживанию систем телемеханики и сетевой инфраструктуры}
\item \textbf{Плановое годовое обслуживание систем телемеханики и гарантированного питания:} устройства телемеханики «Телеканал-М2», СГП (системы гарантированного питания), КТЦ (контрольно-тренировочный цикл) АКБ -- обеспечение бесперебойной работы каналов связи и питания диспетчерского управления.
\item \textbf{Квартальное обслуживание систем сбора и передачи данных:} ССПИ TOPAZ, устройства цифрового контроля и сигнализации (УЦКС-1/2/1M), АРМ телемеханики -- диагностика, профилактика отказов, поддержание готовности каналов передачи телеметрии.
\item \textbf{Внедрение шкафа ССПИ на базе оборудования TOPAZ:} монтаж, сборка и пуско-наладка, настройка перечня сигналов для передачи в ПМЭС (предприятие магистральных электрических сетей). Stack: SNTP, ModBUS/RTU, RS232/485, Ethernet, SSH, FTP, LAN, Visio, Draw.io
\item \textbf{Разработка SCADA-проекта мониторинга бесперебойного гарантированного электропитания ВПМЭС:} проектирование архитектуры сбора данных с ИБП и сопутствующего оборудования, интеграция протоколов сбора телеметрии. Stack: MasterSCADA 3.9, OPC, SNMP, ModBUS/RTU, TCP/IP, Ethernet, Python, CMD, PowerShell, SSH, LAN, Cisco, MOXA
\item \textbf{Миграция автоматизированных рабочих мест на ALT Linux:} самостоятельная установка и настройка более 100 станций в четырёх филиалах предприятия (Череповец, Вологда, Иваново, Кострома). Stack: SMTP, DNS, статическая маршрутизация, FTP, NTFS, FAT32, ExFAT, firewall, LAN
\item \textbf{Актуализация и систематизация сигналов телемеханики:} проверка и внесение в формуляры приёмо-передачи данных параметров активной и реактивной мощности (на линиях и трансформаторах), напряжений на линиях и секциях шин, тока на линиях и выключателях, частоты, дискретных сигналов с выключателей и оповещений об авариях. Stack: PostgreSQL, Visio, Draw.io
\end{joblong}


% \begin{joblong}{Company}{Mar 2019 - Jan 2021}{Designation}
% \item long long line of blah blah that will wrap when the table fills the column width
% \item again, long long line of blah blah that will wrap when the table fills the column width but this time even more long long line of blah blah. again, long long line of blah blah that will wrap when the table fills the column width but this time even more long long line of blah blah
% \end{joblong}



% %Projects
% \section{Projects}

% \begin{tabularx}{\linewidth}{ @{}l r@{} }
% \textbf{Some Project} & \hfill \href{https://some-link.com}{Link to Demo} \\[3.75pt]
% \multicolumn{2}{@{}X@{}}{long long line of blah blah that will wrap when the table fills the column width long long line of blah blah that will wrap when the table fills the column width long long line of blah blah that will wrap when the table fills the column width long long line of blah blah that will wrap when the table fills the column width}  \\
% \end{tabularx}



% %----------------------------------------------------------------------------------------
% %	PUBLICATIONS
% %----------------------------------------------------------------------------------------
% \section{Publications}
% \begin{refsection}[citations.bib]
% \nocite{*}
% \printbibliography[heading=none]
% \end{refsection}



% \vfill
% \center{\footnotesize Last updated: \today}

\end{document}
